\documentclass{article}
\usepackage[utf8]{inputenc}
\usepackage[T1]{fontenc}
\usepackage[french]{babel}
\usepackage{geometry}
\usepackage{fancyhdr}
\linespread{1.15}

% Configuration de la géométrie (marges par défaut)
\geometry{
    a4paper,
    margin=1in,
}

% Configuration des en-têtes et pieds de page
\pagestyle{fancy}
\fancyhf{}
\rhead{MMN1, A4 ESILV}
\lhead{MACHINE LEARNING PROJECT REPORT}
\rfoot{\thepage}
\cfoot{}

% Redéfinition du style de la première page pour qu'elle soit une page de titre simple
\fancypagestyle{firstpage}{
    \fancyhf{}
    \renewcommand{\headrulewidth}{0pt}
    \renewcommand{\footrulewidth}{0pt}
}

% Pour la page de titre, on enlève l'en-tête/pied de page fancy
\thispagestyle{firstpage}

\begin{document}

% --- PAGE 1: PAGE DE TITRE ---
\begin{titlepage}
    \centering
    \vspace*{2cm}
    {\Huge \textbf{MACHINE LEARNING PROJECT REPORT}} \\
    \vspace{1cm}
    {\Large Prediction of the grades of students} \\
    \vfill
    {\Large MMN1, A4 ESILV} \\
    {\Large December 2025} \\
    \vspace{0.5cm}
    {\Large Daphne BARAY, Sofiane BEAUMONT, Max BAUMBERGER} \\
    \vfill
\end{titlepage}

\clearpage
% On remet le style de page fancy après la page de titre
\pagestyle{fancy}

% --- PAGE 2: TABLE DES MATIÈRES ---
\setcounter{page}{2}
\tableofcontents
\thispagestyle{fancy} % S'assurer que la page 2 a bien l'en-tête/pied de page
\clearpage

% --- PAGE 3: INTRODUCTION ---
\section*{Introduction}
\addcontentsline{toc}{section}{Introduction}
\label{sec:intro}
\setcounter{page}{3}
% Le contenu de l'introduction irait ici
\clearpage

% --- PAGE 4: PARTIE 1 ---
\section*{Partie 1}
\addcontentsline{toc}{section}{Partie 1}
\label{sec:partie1}
\setcounter{page}{4}
\subsection*{Sous-partie 1}
\addcontentsline{toc}{subsection}{Sous-partie 1}
\label{sec:partie1_s1}
% Contenu de la Sous-partie 1
\subsection*{Sous-partie 2}
\addcontentsline{toc}{subsection}{Sous-partie 2}
\label{sec:partie1_s2}
% Contenu de la Sous-partie 2
\clearpage

% --- PAGE 5: PARTIE 2 ---
\section*{Partie 2}
\addcontentsline{toc}{section}{Partie 2}
\label{sec:partie2}
\setcounter{page}{5}
\subsection*{Sous-partie 1}
\addcontentsline{toc}{subsection}{Sous-partie 1}
\label{sec:partie2_s1}
% Contenu de la Sous-partie 1
\subsection*{Sous-partie 2}
\addcontentsline{toc}{subsection}{Sous-partie 2}
\label{sec:partie2_s2}
% Contenu de la Sous-partie 2
\clearpage

% --- PAGE 6: PARTIE 3 ---
\section*{Partie 3}
\addcontentsline{toc}{section}{Partie 3}
\label{sec:partie3}
\setcounter{page}{6}
\subsection*{Sous-partie 1}
\addcontentsline{toc}{subsection}{Sous-partie 1}
\label{sec:partie3_s1}
% Contenu de la Sous-partie 1
\subsection*{Sous-partie 2}
\addcontentsline{toc}{subsection}{Sous-partie 2}
\label{sec:partie3_s2}
% Contenu de la Sous-partie 2
\clearpage

% --- PAGE 7: RÉSULTATS ---
\section*{Résultats}
\addcontentsline{toc}{section}{Résultats}
\label{sec:resultats}
\setcounter{page}{7}
% Le contenu des résultats irait ici
\clearpage

% --- PAGE 8: CONCLUSION ---
\section*{Conclusion}
\addcontentsline{toc}{section}{Conclusion}
\label{sec:conclusion}
\setcounter{page}{8}
% Le contenu de la conclusion irait ici

\end{document}